\input{../tex-inputs/latex-korrekturansicht-vorspann}

\section[Arthur Schnitzler an Theodor Herzl, 5. 7. 1894]{L03908 Arthur Schnitzler an Theodor Herzl, 5. 7. 1894}
\nopagebreak\mylabel{L03908v}
\rehead{ }\normalsize\beginnumbering\briefempfaengerindex{Herzl, Theodor@\textsc{Herzl, Theodor}!zzzSchnitzler, Arthur@\emph{von Arthur Schnitzler}!1894-07-051@{5. 7. 1894}|(be}
\toendnotes[C]{\smallbreak\pagebreak[2]}
\correspDesc{Versand  durch Arthur Schnitzler am 5. 7. 1894 in Wien
\newline{}Erhalt  durch Theodor Herzl im Zeitraum [5. 7. 1894
                  – 8. 7. 1894?] in Wien}\toendnotes[C]{\smallbreak}
\Standort{Jerusalem, Central Zionist Archives, H1:1924-13.}
\physDesc{Brief, 2 Blätter, 8 Seiten, 2649 Zeichen
\newline{}Handschrift: schwarze Tinte, deutsche Kurrent
\newline{}Ordnung: mit Bleistift von unbekannter Hand innerhalb das Konvoluts paginiert:
                                    »39«–»41« }
\buchAbdrucke{\weitereDrucke{Arthur Schnitzler: \emph{Briefe 1875–1912}. Herausgegeben von Therese Nickl und Heinrich Schnitzler. Frankfurt am Main: \emph{S. Fischer} 1981, S. 228–229.} }\toendnotes[C]{\smallbreak}
\pstart
           \raggedleft{}{\pb}\textsc{\textcolor{pink}{IX. Frankgasse 1}\oindex{Wien@\textbf{Wien}!IX., Alsergrund@\textbf{IX., Alsergrund}!Frankgasse 1@\textbf{Frankgasse 1}, \emph{Wohngebäude}|pw}{}\ledrightnote{\textcolor{pink}{Frankgasse 1}}.}\pend
           
\pstart{}Verehrter Freund,\pend\vspace{0.5em}
\pstart
           Ihre \label{K_L03908-1v}\edtext{freundlichen Worte}{\lemma{\textnormal{\emph{freundlichen Worte}}}\Cendnote{\textnormal{Theodor Herzl an Arthur Schnitzler, 3. 7. 1894.}}}\label{K_L03908-1} haben meine \textcolor{blue}{Mutter}\pwindex{Schnitzler, Louise 8.\,7.\,1840 Kőszeg – 9.\,9.\,1911 Wien@\textsc{Schnitzler, Louise} (8.\,7.\,1840 Kőszeg – 9.\,9.\,1911 Wien)|pwv}{}\ledrightnote{{$\rightarrow$}\emph{\textcolor{blue}{Louise Schnitzler}}} und meinen \textcolor{blue}{Bruder}\pwindex{Schnitzler, Julius 13.\,7.\,1865 Wien – 29.\,6.\,1939 ebd.@\textsc{Schnitzler, Julius} (13.\,7.\,1865 Wien – 29.\,6.\,1939 ebd.), \emph{Chirurg}|pwv}{}\ledrightnote{{$\rightarrow$}\emph{\textcolor{blue}{Julius Schnitzler}}} ſehr erfreut, und ich
               danke Ihnen in ihrem und meinem Namen aufs wärmſte. –\pend
           
\pstart
           Daſs ich Ihnen nichts von meinen Sachen ſchicke, nach denen Sie ſich in ſo
               liebenswürdiger Weiſe erkundigen, liegt wirklich weniger an mir als an den Verlegern,
               die ſich noch {\pb}immer ſehr lang bitten laſſen, bevor ſie was
               von mir drucken. Nun, im Herbſt erſcheint eine \textcolor{green}{Novelle}\pwindex{Schnitzler, Arthur 15. 5. 1862 Wien – 21. 10. 1931 ebd.@\textsc{Schnitzler, Arthur} (15. 5. 1862 Wien – 21. 10. 1931 ebd.), \emph{Schriftsteller, Mediziner}!Sterben. Novelle@\strich\emph{Sterben. Novelle}|pwv}{}\ledrightnote{{$\rightarrow$}\emph{\textcolor{green}{Sterben. Novelle}}} von mir bei \textsc{\textcolor{brown}{Fischer}\orgindex{S. Fischer Verlag@S. Fischer Verlag|pw}{}\ledrightnote{\textcolor{brown}{S. Fischer Verlag}}}, und ich will mir alle Mühe geben, anderes, das nun ſchon fertig im Pult liegt,
               raſcher in die Oeffentlichkeit zu befördern, als es mir bisher zu gelingen pflegte.
               Im ganzen darf ich ſagen, dß ich in den letzten Monaten nicht {\pb}ſehr nachläßig war, daſs mir mancherlei einfällt und dß ich
               zuweilen die Empfindung habe, daſs ich manches von dieſem Mancherlei werde zu gutem
               Ende führen können. –\pend
           
\pstart
           Ich zweifle nicht, dß mein Freund \textcolor{blue}{Paul}\pwindex{Goldmann, Paul 31.\,1.\,1865 Breslau – 25.\,9.\,1935 Wien@\textsc{Goldmann, Paul} (31.\,1.\,1865 Breslau – 25.\,9.\,1935 Wien), \emph{Schriftsteller, Journalist}|pw}{}\ledrightnote{\textcolor{blue}{Paul Goldmann}} Ihnen
               meine Grüße \introOben{}an Sie,\introOben{} die ich den \label{K_L03908-2v}\edtext{Briefen an ihn}{\lemma{\textnormal{\emph{Briefen an ihn}}}\Cendnote{\textnormal{Diese Briefe sind nahezu
               vollständig verschollen.}}}\label{K_L03908-2} häufig beifüge, regelmäßig beſtellt,
               und {\pb}Ihnen auch manchmal ſagt – was ſich mündlich und durch
               einen Dritten beſſer ſagen läßt als in einem Brief, wo es einen süßlich faden
               Beigeſchmack von Höflichkeit oder gar Förmlichkeit beko{\geminationm}t – nemlich daſs ich das wenige, was mir von Ihnen zugänglich iſt\textcolor{gray}{,} ſtets mit
               wahrhaftem Genuße leſe. Beſonders im Laufe des letzten Jahres haben Sie einige kleine
               Kunſtwerke {\pb}von Feu{[}i{]}lletons
               geſchaffen, die nicht mit den Zeitungen selbst verwehen dürften. Sie wiſſen das
               ſelbſt und man darf es Ihnen wohl ſo unbefangen ins Geſicht ſagen wie eine Grobheit.
               – Und die Bühne? Iſt Ihre Luſt zum Dramatiſchen gänzlich durch den Ekel erſtickt
               worden? Wie oft hab’ ich in dieſem Winter an Ihre \label{K_L03908-3v}\edtext{ſchönen und wahren Worte}{\lemma{\textnormal{\emph{ſchönen und wahren Worte}}}\Cendnote{\textnormal{Theodor Herzl an Arthur Schnitzler, 29. 7. 1892.}}}\label{K_L03908-3} denken müſſen, die Sie mir
               lang vor der Aufführung meines »\textcolor{green}{Märchens}\pwindex{Schnitzler, Arthur 15. 5. 1862 Wien – 21. 10. 1931 ebd.@\textsc{Schnitzler, Arthur} (15. 5. 1862 Wien – 21. 10. 1931 ebd.), \emph{Schriftsteller, Mediziner}!Märchen. Schauspiel in drei Aufzügen@\strich\emph{Das Märchen. Schauspiel in drei Aufzügen}|pw}{}\ledrightnote{\textcolor{green}{Das Märchen. Schauspiel in drei Aufzügen}}«
               geſchrieben haben. {\pb}Ich habe von allem zu koſten beko{\geminationm}en, was die Aufführung eines Stückes verletzendes
               bringen kann: wie irgend einer ka{\geminationn} ich mitreden, wenn
               von der Albernheit des Direktors, der Verlegenheit der Komödianten und der vergnügten
               Gefälligkeit der Recenſenten geſprochen wird, – wobei ich vom Publikum gänzlich
               ſchweigen will, das albern, verlogenen und gefällig iſt. – Es iſt nicht anzunehmen,
               daſs ich anders reden würde, wenn {\pb}ich zufällig einen
               Erfolg gehabt hätte, nur ſetzte ich hinzu: Trotzdem{\dotsfour}{ }\textsc{etc.} – und es klänge großartiger. –\pend
           
\pstart
           – Wenn Sie in \textcolor{pink}{Auſſee}\oindex{Bad Aussee@\textbf{Bad Aussee}, \emph{Hauptstadt}|pw}{}\ledrightnote{\textcolor{pink}{Bad Aussee}} ſein werden, ſo hoffe ich
               die Freude zu haben Sie zu ſehen, da ich im Auguſt meine \textcolor{blue}{Mama}\pwindex{Schnitzler, Louise 8.\,7.\,1840 Kőszeg – 9.\,9.\,1911 Wien@\textsc{Schnitzler, Louise} (8.\,7.\,1840 Kőszeg – 9.\,9.\,1911 Wien)|pwv}{}\ledrightnote{{$\rightarrow$}\emph{\textcolor{blue}{Louise Schnitzler}}} in \textcolor{pink}{Iſchl}\oindex{Bad Ischl@\textbf{Bad Ischl}|pw}{}\ledrightnote{\textcolor{pink}{Bad Ischl}} beſuchen werde. Vielleicht laſſen Sie aber bis dahin noch ein
               freundliches Wort von ſich hören. Haben Sie die Güte mich Ihrer w. Frau \textcolor{blue}{Gemahlin}\pwindex{Herzl, Julie 1.\,2.\,1868 Budapest – 10.\,11.\,1907 Wien@\textsc{Herzl, Julie} (1.\,2.\,1868 Budapest – 10.\,11.\,1907 Wien)|pwv}{}\ledrightnote{{$\rightarrow$}\emph{\textcolor{blue}{Julie Herzl}}} beſtens zu empfehlen,
               und ſeien Sie, mein\pend
           
\pstart
           {\pb}lieber Freund, aufs herzlichſte{\\[\baselineskip]}gegrüßt. Ihr
                  \spacefill\mbox{ArthurSchnitzler}\pend
           \leftskip=0em{}
\pstart
           \textcolor{pink}{Wien}\oindex{Wien@\textbf{Wien}, \emph{Verwaltungsgebiet}|pw}{}\ledrightnote{\textcolor{pink}{Wien}}, 5. 7. 94\pend
           \selectlanguage{ngerman}\endnumbering\briefempfaengerindex{Herzl, Theodor@\textsc{Herzl, Theodor}!zzzSchnitzler, Arthur@\emph{von Arthur Schnitzler}!1894-07-051@{5. 7. 1894}|)be}\mylabel{L03908h}  \newcommand{\dateiname}{L03908}\newcommand{\titel}{Arthur Schnitzler an Theodor Herzl, 5. 7. 1894}\newcommand{\editorInnen}{Herausgegeben von Jahnke, SelmaMüller, Martin Anton}\input{../tex-inputs/latex-korrekturansicht-abspann}
